\section{Conclusion}
\label{sec:conclusion}

This paper has revisited the canonical three-equation New Keynesian
model and characterised, in closed form, the dynamic response to a 25
basis-point contractionary monetary policy shock. Three findings stand
out. On impact, output falls by $0.26\%$, annualised inflation falls by
$0.35$ percentage points, and the annualised real interest rate rises
by $0.52$ percentage points. The three responses converge to the
steady state within approximately ten quarters and are robust to
plausible variation in the Taylor-rule reaction coefficients.

The exercise abstracts from several features that matter empirically.
Capital accumulation, sticky wages, habit formation in consumption,
and a richer set of shocks would enrich the propagation mechanism and
in some cases substantially alter the quantitative magnitudes
\citep[cf.][]{SmetsWouters2007}. Financial frictions would change the
way monetary policy is transmitted through credit markets. Quantifying
these extensions within the same transparent solution framework
adopted here is a natural next step. The transparent quantitative
benchmark documented in this paper provides a baseline against which
the marginal contribution of each extension can be assessed.
